%auto-ignore
    \begin{figure}[ht]
        \centering
        \begin{tikzpicture}[
          inner/.style={draw,fill=blue!5,thick,inner sep=3pt,minimum width=8em},
          outer/.style={draw=gray,dashed,fill=green!1,thick,inner sep=5pt}
          ]

            \node[draw,
                inner,
                minimum width=2cm,
                minimum height=1cm](param) at (0,0) {Parameter estimation};

            \node[,
                above= of param,
                minimum width=2cm,
                minimum height=1cm](train1) {Training data \((X,\bm{y})\)};

            \node[draw,
                inner,
                right=of param,
                minimum width=2cm,
                minimum height=1cm](pred) {Prediction};

            \node[,
                above= of pred,
                minimum width=2cm,
                minimum height=1cm](train2) {Training data \((X,\bm{y})\)};

            \node[,
                below= of pred,
                minimum width=2cm,
                minimum height=1cm](target) {Target point \(\bm{x}^*\)};

            \node[draw,
                    right= of pred,
                    minimum width=2cm,
                    minimum height=1cm] (dist) {Distribution of \(y^*\given X,\bm{y}\) at
                    \(\bm{x}^*\)};
                    
            %arrows
            \draw[-latex] (param) edge [below, "$\hat{\bm{\theta}}$"] (pred);

            \draw[-latex] (pred) edge (dist);
           
            \draw[-latex] (train1) edge (param);

            \draw[-latex] (train2) edge (pred);

            \draw[-latex] (target) edge (pred);

            \begin{pgfonlayer}{background}
                \node[outer,fit=(param) (pred)] (joint) {};
            \end{pgfonlayer}

    \end{tikzpicture}
    \caption{Flowchart of the GP regression procedure. The dashed box indicates the usual approach of combining the parameter estimation and prediction tasks under one strategy.}
    \label{fig:flowchart}
\end{figure}